\newcommand{\titreA}{Rapport intermédiaire}
\newcommand{\titreB}{Projet industriel}
\documentclass[12pt]{article}
\usepackage[utf8]{inputenc}
\usepackage[francais]{babel}
\usepackage[T1]{fontenc}
\usepackage{lmodern, marvosym, geometry, graphicx, multicol, lastpage, tikz, listings}
\usepackage[hyperindex=true, colorlinks=true, breaklinks=true, linkcolor=blue]{hyperref}
\usepackage{fancyhdr, verbatim, alltt, pdfpages}

\geometry{hmargin=1.5cm, vmargin=2cm}
\addtolength{\parskip}{10pt}
\pagestyle{fancy}

\definecolor{lightgray}{gray}{0.9}
\definecolor{darkgray}{gray}{0.4}
\newcommand{\hlc}[1]{\color{purple}{\textbf{#1}}}
\lstset{language=,
    basicstyle=\sffamily\footnotesize,
    xleftmargin=20pt,
    xrightmargin=20pt,
    numbers=left,
    stepnumber=1, % le pas des numeros de ligne
    numbersep=10pt,
    breaklines=true,
    tabsize=4,
    %frame=single,
    showspaces=false,
    showstringspaces=false
    inputencoding=utf8, %put1 d'utf8 dans listing
    extendedchars=true, %put1 d'utf8 dans listing
    literate={à}{{\`a}}1 {é}{{\'e}}1 {è}{{\`e}}1 {ê}{{\^e}}1 {â}{{\^a}}1 {î}{{\^i}}1 {ê}{{\^e}}1 {É}{{\'E}}1 {À}{{\`A}}1 {«}{{\og}}1 {»}{{\fg{}}}1 {ô}{{\^o}}1 {ù}{{\`u}}1 {û}{{\^u}}1 {ç}{{\c{c}}}1 {Ç}{{\c{C}}}1 {--}{{-\,-}}1 {-}{{-}}1 {*}{{*}}1 {Switch(config)\#}{{\textbf{Switch(config)\# }}}1 {Switch(config-if)\#}{{\textbf{Switch(config-if)\# }}}1 {Switch(config-line)\#}{{\textbf{Switch(config-line)\# }}}1 {Switch>}{{\textbf{Switch> }}}1 {Switch\#}{{\textbf{Switch\# }}}1, %put1 d'utf8 dans listing
    columns=fullflexible, % suppression des espaces autour des deux-points
    escapechar=§, % permet d'insérer du §code latex§
    backgroundcolor=\color{lightgray},
    alsoletter={.,1,2,3,4,5,6,7,8,9,0},
    keywords={path_to_certs, user_login, user_password, nom_entreprise, pass_radius_sql, utilisateur1, pass_utilisateur1, secret_radius, 192.168.1.2, 192.168.1.10, 255.255.255.0, fastethernet0/1, vlan_id, pass_enable, pass_enable_secours, utilisateur_secours, pass_utilisateur_secours, pass_root_sql, snack, 192.168.0.254, commutateur1, client1, Nancy, BHConsulting, France, FR, nom_serveur_radius, dossier_certs, dossier_certs_utilisateur, eth0 },
    keywordstyle=\hlc,
    morecomment=[l]{\#}, % commentaires shell en ligne
    commentstyle=\itshape\color{darkgray},
}

\renewcommand{\headrulewidth}{1pt}
\lhead{\textbf{\titreA{} (\titreB)}}
\rhead{\emph{BOUGET / GUÉPIN / PINHÈDE / VAUBOURG}}
\lfoot{TELECOM Nancy - PI}
\cfoot{\thepage{} / \pageref{LastPage}}
\rfoot{2012-2013}

\begin{document}
\thispagestyle{empty}

\begin{multicols}{2}
{\large
	\begin{flushleft}
		\noindent{}\textbf{Nicolas BOUGET}\\
		\Letter~nicolas.bouget@esial.net\\
		3A TRS1\\~

		\noindent{}\textbf{Julien GUÉPIN}\\
		\Letter~julien.guepin@esial.net\\
		3A IL\\
	\end{flushleft}

	\begin{flushright}
		\noindent{}\textbf{Marc PINHÈDE}\\
		\Letter~marc.pinhede@esial.net\\
		3A LE\\~

		\noindent{}\textbf{Julien VAUBOURG}\\
		\Letter~julien@vaubourg.com\\
		3A TRS2\\
	\end{flushright}
}
\end{multicols}

\vspace{0.5cm}

\begin{center}
	{\Huge\textbf{\titreA}}

	\vspace{1cm}

	{\huge\emph{\titreB}}

	\vspace{1cm}

	\begin{flushleft}
		{\large
		\hspace{3.2cm}
		\textbf{Société~:} B.H. Consulting\\
		\hspace{3.2cm}
		\textbf{Intervenant industriel~:} Guillaume ROCHE\\
		\hspace{3.2cm}
		\textbf{Intervenant universitaire~:} Jean-François SCHEID
		}
	\end{flushleft}

	\vspace{1cm}
	{\large Le \today}

	\vspace{1.5cm}

	\includegraphics[width=140pt]{img/BHConsulting.jpg}

	\vspace{1.5cm}

	\includegraphics[width=190pt]{img/ul.png}
	\hspace{3.5cm}
	\includegraphics[width=140pt]{img/telecom-nancy.jpg}
\end{center}
\newpage

\thispagestyle{empty}
\tableofcontents
\newpage



\section{Introduction}
\section{B.H. Consulting}
\section{Projet}
\subsection{Contexte}

La sensibilisation croissante aux notions de sécurisation des réseaux des entreprises a conduit B.H. Consulting à devoir répondre à de nouveaux types d'exigences.

Les réseaux d'une entreprise sont en effet devenus un point critique de toutes les nouvelles sociétés, qui sont régulièrement amenées à perdre des heures considérables de travail dès lors que ceux-ci ne sont plus accessibles. Au-delà de l'aspect fonctionnel, les réseaux sont une porte vers les serveurs les plus sensibles des administrations, qui y stockent les informations les plus confidentielles.

Qu'ils soient filaires ou non, ces réseaux sont accessibles dès lors qu'on est dans les locaux des entreprises~: prises murales, émissions wifi, etc. N'importe quel individu ayant la possibilité de se relier physiquement aux réseaux, il faut mettre en place un mécanisme de contrôle logique pour restreindre ses accès. Pour autant, ce mécanisme ne doit pas être lourd à administrer, ce qui impliquerait rapidement un laxisme et une confusion dans les règles établies. Définir strictement les autorisations en fonction des points d'accès physiques, pour ensuite établir une carte des autorisations de l'ensemble des bâtiments n'est pas non plus la bonne solution, puisqu'il sera difficile de garantir les contrôles d'accès à ces prises. L'utilisation du réseau dépendra également de  l'emplacement géographique de l'utilisateur, à l'encontre de toute logique vis-à-vis de la tendance actuelle au nomadisme constaté dans les entreprises.

Pour répondre à toutes ces exigences, critiques pour une entreprise qui souhaite garantir l'intégrité de son réseau et la confidentialité de ses données tout en sauvegardant la liberté de mouvement de ses employés, B.H. Consulting est régulièrement contraint d'exploiter les nouvelles technologies les plus adéquates. La première problématique de ce projet industriel sera donc de permettre à B.H. Consulting d'avoir à disposition la documentation nécessaire pour déployer l'une des solutions existantes, facilement et rapidement.

Au delà du déploiement, B.H. Consulting souhaite offrir une souplesse d'administration exemplaire de son système, en permettant aussi bien à ses clients que ses propres employés de modifier les permissions des réseaux. Aucune des solutions d'interface existantes ne répondant à ce besoin actuellement, il s'agit du second objectif de ce projet industriel. L'interface devra devra donc être accessible, ergonomique et complète. Des fonctionnalités d'administration poussées comme la gestion d'un historique, d'une console virtuelle et la gestion des équipements actifs la compléteront durant ce projet.

Les objectifs de ce projet industriel émanent directement des besoins des clients, auxquels B.H. Consulting doit faire face.

\subsection{Objectifs}

Plusieurs objectifs peuvent être tirés des problématiques sus-évoquées~:

\begin{enumerate}
\item Définir la technologie qui permettra de définir des accès aux réseaux de façon dynamique, souple et sécurisée.
\item Former l'équipe du projet industriel sur cette technologie, à l'aide de documentations et d'expérimentations concrétes sur les différents aspects sous-jacents.
\item Écrire une documentation sur le sujet qui sera restreinte et utile à B.H. Consulting dans le cadre de ses besoins.
\item Concevoir une solution logicielle permettant le déploiement rapide et efficace de la technologie sur les réseaux des clients.
\end{enumerate}

Une seconde liste d'objectifs de dégage rapidement, et de façon relativement indépendante de la première~:

\begin{enumerate}
\item Concevoir une interface permettant la configuration des accès au réseau, avec des profils d'utilisateurs définis par des ensembles de permissions.
\item Ajouter des fonctionnalités supplémentaires à l'interface, comme la console ou la gestion des équipements actifs.
\item Intégrer l'interface dans la solution logicielle proposée lors de la première étape, pour le déploiement de la solution chez les clients.
\end{enumerate}

Afin de répondre aux attentes de l'entreprise de façon optimale, en ayant l'assurance d'être exhaustif et de soulever toute ambiguïté, les objectifs ont nécessité d'être contractualisés. Les documents qui en résultent ont été rédigés en partenariat avec les responsables de l'entreprise, en respectant les étapes classiques d'une gestion de projet.

\section{Gestion de projet}
\subsection{Constitution de l'équipe}

L'équipe du projet est constituée exclusivement d'élèves-ingénieurs de l'école TELECOM Nancy, en dernière année de formation. Après avoir remporté la confiance de l'entreprise pour se faire attribuer ce travail, le chef de projet a confirmé l'équipe suivante~:

\begin{description}
\item[Nicolas BOUGET] Ayant intégré la spécialisation \textit{Télécoms, Réseaux et Services} (TRS) dès la seconde moitié de la deuxième année de son cursus ingénieur, Nicolas a une sensibilité particulière pour le réseau et l'administration de services. Avec une expérience forte de conception de tests pour un simulateur d'avion, acquise durant son stage de deuxième année, Nicolas dispose d'une rigueur et d'une méthodologie reconnues, avec une sensibilisation aux problématiques du développement de gros projets. Co-responsable du parc informatique de la convention de culture japonaise Anim'Est lors de sa dernière édition, il est également capable de gérer des infrastructures imposantes en supportant un stress important et une nécessité de qualité de service exemplaire. C'est donc naturellement qu'il a été désigné comme chef de projet, dans le cadre de ce travail.\\
\item[Julien GUÉPIN] Julien a choisi la spécialisation \textit{Ingénieurie Logicielle} (IL) lors de sa seconde année à TELECOM Nancy. Véritable passionné de développement web, Julien apporte des compétences extrêmement fortes dans un projet qui nécessite une expertise sérieuse pour aboutir à une solution logicielle munie d'une interface web fonctionnelle, ergonomique et efficace. Avec de nombreuses expériences dans le domaine, il a été capable de prouver ses capacités de nombreuses fois dans le passé, autant au travers de ses stages que ses projets personnels. Actif associativement, Julien s'est notamment investi dans un projet humanitaire au Pérou, développant ainsi des capacités évidentes de travail en équipe et d'organisation personnelle pour conduire un projet complexe vers une réussite reconnue. Particulièrement intéressé par la seconde partie du projet, il est un élément clé de sa réussite.\\
\item[Marc PINHÈDE] C'est une troisième spécialité offerte par l'école que Marc propose en intégrant l'équipe, puisqu'il a choisi dès l'année précédente l'option \textit{Logiciels Embarqués} (LE). Sa passion pour le domaine, ainsi que la rigueur imposée par les environnements embarqués font de lui une ressource particulièrement minutieuse et attachée à la perfection, qui conduit à la performance et l'excellence. Marc a notamment eu l'occasion de prouver ses compétences lors de son stage de seconde année dans un centre de recherche, en y implémentant un compilateur. Avec des capacités d'intégration et de travail en équipe reconnues à chacune de ses implications dans la convention Anim'Est, Marc est un membre essentiel de l'équipe, qui saura mettre à profit ses compétences dans les domaines couverts par ce projet, qui l'intéressent tout autant.\\
\item[Julien VAUBOURG] Également issue de la spécialisation TRS, Julien est un passionné des réseaux qui dispose en plus d'une forte expérience dans le développement. Une double-compétence qu'il a su prouver à plusieurs reprises au travers de ses nombreux stages et projets personnels. Particulièrement impliqué associativement, notamment dans l'administration d'un fournisseur d'accès à Internet participatif, Julien a l'occasion d'intégrer très régulièrement des équipes, tout en étant force de proposition. Également acteur de la dernière édition de la convention Anim'Est auprès de Nicolas et Marc, il intègre naturellement l'équipe avec une capacité de vision d'ensemble du projet.
\end{description}

Il s'agit donc d'une équipe particulièrement pluridisciplinaire, qui saura composer avec ses membres pour aggréger les compétences de chacun afin d'obtenir le meilleur d'eux-mêmes. Tous les membres ayant déjà eu l'occasion de travailler ensemble, l'équipe est dynamique et prête à travailler dès le lancement du projet.

\subsection{Note de cadrage}

Suite à une réunion avec l'intervenant industriel, la note de cadrage a été rédigée. L'objectif de cette note de cadrage était de fournir une base qui nous a permise de discuter ensemble des objectifs et contraintes du projet, de façon plus contractuelle qu'une simple conversation informelle.

Après avoir échangés sur cette base de travail, nous sommes arrivés à une note de cadrage satisfaisant l'encadrant et l'équipe du projet.

La note de cadrage rappelle la problématique et les bénéfices du projet pour l'entreprise, avec un effort de vulgarisation particulier pour la rendre accessible à un non-technicien. Elle détaille précisément le fonctionnement de la technologie retenue en accord avec l'intervenant, et l'objectif de son utilisation pour l'entreprise. 

Un contexte est ensuite proposé, en rappelant le statut particulier du projet industriel et les enjeux qui sont propres aux membres de l'équipe. Le délai ainsi que l'absence de budget sont évoqués, avec le matériel qui est disposition au lancment du projet.

Les objectifs sont complétés par un chapitre sur le niveau de qualité attendu, avec une priorisation claire des principales tâches à traiter. Une section risques rappelle les points sensibles du projet.

La liste des livrables y est également définie~:
\begin{description}
\item[19 octobre 2012 :] Remise de la note de cadrage.
\item[5 novembre 2012 :] Signature de la note de cadrage.
\item[12 novembre 2012 :] Remise du cahier des charges.
\item[19 novembre 2012 :] Signature du cahier des charges.
\item[20 décembre 2012 :] Soutenance intermédiaire.
\item[7/8 février 2013 :] Audit du projet.
\item[14 mars 2013 :] Remise du projet avec les documentations.
\item[20 mars 2013 :] Remise du rapport.
\item[21 mars 2013 :] Soutenance finale.
\end{description}

Comme convenu, la note de cadrage a été signée par le chef de projet et l'intervenant industriel. C'est en partant de son contenu fiable que le cahier des charges a été établit.

\subsection{Cahier des charges}

Le cahier des charges a été fourni par l'entreprise.

Il rappelle brièvement le matériel qui est à notre disposition durant toute la durée de ce projet. La description des tâches, organisée en liste à puces, est découpée en deux parties. La première concerne les authentifications Radius, tandis que la seconde concerne la gestion des configurations Cisco des équipements actifs du réseau.

Les différentes tâches, ainsi que la description des différents concepts techniques, sont expliqués dans la suite de ce document.

\subsection{Organisation}
\subsubsection{Planning}
\subsubsection{Outils de gestion}

Pour une organisation optimale du travail fourni, l'équipe a recouru à une solution logicielle de gestion de projet.

L'équipe a défini ensemble lors des premières réunions les principales tâches qui constituaient ce projet, en dégageant notamment les deux étapes majeures. Le nombre de jour-hommes a été évalué pour chacune d'entre elles, permettant ainsi d'aboutir à une diagramme de Gantt qui a précise l'enchainement logique de ces tâches.

Toutes les grandes tâches qui pouvaient être parallélisées et qui débutaient dès le lancement du projet ont été revues en détail, pour obtenir une vision plus précise de l'organisation de cette période du planning. Au fur et à mesure que le projet avançait, les sous-tâches ont été définies en réunion et se sont ajoutées aux étapes principales, permettant à l'ensemble de l'équipe d'avoir une vision globale pour le long terme et une vision claire pour le court terme, en prenant en compte les contraintes rencontrées qui n'étaient initialement pas prévues.

L'outil de gestion permettant d'attribuer les tâches aux différents membres de l'équipe, aucune d'entre elles qui aurait pu mettre l'avancement du projet en péril n'a été laissée de côté. Enfin, les membres de' léquipe ont eu l'occasion d'ajouter des feuilles de temps chaque fois qu'ils ont travaillé sur l'une de leurs tâches en particulier, permettant ainsi de suivre progressivement le bon avancement de celles-ci.

Chaque réunion a donné lieu à un rapport, transmis à l'ensemble des membres de l'équipe ainsi qu'à l'intervenant industriel et l'encadrant universitaire.

\subsubsection{Attributions et réunions}

L'attribution des tâches aux membres du projet s'est faite sur la base du volontariat. Le groupe étant cohérent et dynamique, aucune tâche n'a été affectée de force, et cette méthode a fonctionné correctement tout au long du projet, permettant à chacun d'avancer sans contrariétés.

Une réunion de travail a été planifiée toutes les semaines, en présence de l'ensemble des membres de l'équipe et de l'intervenant industriel qui s'est montré disponible tout au long du projet. Une réunion mensuelle a eu lieu avec l'encadrant universitaire, lui permettant d'obtenir une interactivité régulièrement complémentaire avec les rapports des réunions de travail reçus toutes les semaines.

\section{Choix effectués}
\subsection{Protocoles et implémentations}
\subsubsection{Contrôle des accès réseau}

Lors des premières réunions avec l'intervant industriel, il a été convenu d'utiliser la technologie 802.1x pour répondre aux exigences du projet. Cette solution correspond particulièrement bien avec les objectifs, et est déjà connu au sein de l'entreprise.

Le 802.1x est un standard du réseau, largement utilisé dans les réseaux d'entreprise. Il s'agit d'un protocole implémenté au niveau des équipements actifs du réseau : les commutateurs et les routeurs. Lorsqu'un poste client tente de se brancher physiquement à la prise murale de son bureau, celle-ci est désactivée par défaut. Une fois l'ordinateur relié, celui-ci prend connaissance de la présence de la technologie 802.1x sur le réseau, et interroge le commutateur auquel il accède. Pour ce faire, il lui envoie ce qui est nécessaire à son autthentification. Selon le type d'authentification reconnu sur le réseau, il y a plusieurs possibilités~:

\begin{description}
\item[Par adresse MAC~:] Il s'agit de la méthode la plus simple à mettre en oeuvre, et à utiliser pour le poste client, puisqu'elle est totalement transparente. L'adresse MAC représente l'adresse physique de l'interface réseau du poste client, unique dans le monde. Elle suffit donc au commutateur pour identifier celui qui tente de rejoindre le réseau, et l'autoriser ou non à y accéder. Le principal défaut de cette solution est qu'il est très facilement d'usurper l'adresse d'un autre poste client, et de se faire passer pour lui, bénéficiant ainsi automatiquement de ses autorisations.
\item[Par mot de passe~:] En plus, ou non, de l'authentification par adresse MAC, l'authentification peut nécessiter un couple identifiant/mot de passe. Cette méthode résout le problème d'usurpation d'adresse, puisque le mot de passe assure qu'au moins une information est suffisamment secrète pour se prévenir des vols d'identité. Le défaut majeur de cette possibilité est qu'elle n'est plus transparente pour l'utilisateur, qui devra saisir son identifiant et son mot de passe pour rejoindre le réseau. D'un point de vue de la sécurité, le mot de passe saisi transite en clair sur le réseau, permettant ainsi à un éventuel attaquant de le récupérer pour une future usurpation d'identité.
\item[Par certificat~:] Une dernière solution consiste à utiliser un système de certificats. Il en existe de plusieurs types, mais leur point commun à tous est de permettre aux postes clients de s'authentifier de façon totalement transparente, tout en assurant un chiffrement de la communication de bout en bout. Son défaut majeur, c'est que s'il est transparent à l'utilisation, il nécessite une première installation sur les postes clients concernés.
\end{description}

Après discussions entre membres de l'équipe sur les avantages et les inconvénients des différentes solutions, un exposé des problématiques a été présenté à l'intervenant industriel lors d'une réunion hebdomadaire. Chacune de ces solutions ayant ses avantages, l'intervenant a souhaité que toutes ces solutions soient implémentés dans la solution logicielle finale. Malgré sa faible fiabilité, l'authentification par adresses mac est parfaitement adapté aux périphériques commes les imprimantes, l'authentification par mot de passe est utile lorsqu'il s'agit de connexions ponctuels d'utilisateurs invités, et l'authentification par certificats est nécessaire pour tous les permanents de l'entreprise qui bénéficient de permissions importantes sur le réseau.

Lorsque le poste client envoie les informations permettant son authentification au commutateur qui se situe à l'autre bout de la prise murale, ce dernier transmet automatiquement les informations à un serveur chargé de l'authentification et de la mise en place des autorisations. Si l'authentification est bonne, le serveur indique au commutateur l'action qu'il doit entreprendre dans le cadre du protocole 802.1x. Soit le port reste fermé, et le poste client n'a pas accès au réseau, soit il s'ouvre et le réseau est accessible. Selon le type d'utilisateur reconnu, le serveur peut spécifier un numéro de réseau virtuel (VLAN) en particulier. Cette particularité aura pour effet de placer automatiquement le poste client dans un réseau qui bénéficie plus ou moins de privilèges d'accès.

Cette technologie présente une second avantage pour l'entreprise~: en cas de remplacement du commutateur, la configuration du nouveau matériel sera rapide à remettre en place, dans la mesure où toutes les informations sont stockées au niveau du serveur d'authentification et les ports sont entièrement dynamiques.

Les serveurs d'authentification les plus connu dans le domaine des réseaux et en particulier du 802.1x sont ceux qui sont basés sur le protocole Radius. Il en existe plusieurs implémentations, parmies lesquelles un choix a dû être fait.

\subsection{Serveur d'authentification}

Le marché dispose de plusieurs dizaines d'implémentations différentes du protocole Radius. L'entreprise souhaite se restreindre aux solutions libres, le choix se réduit automatiquement à une dizaine de solutions différentes.

Après avoir éliminé les solutions trop anecdotiques, qui ne bénéficient ni d'un développement suivi ni d'une vraie communauté lui permettant d'assurer son avenir, trois implémentations ont été retenues~:

\begin{itemize}
\item FreeRadius
\item GnuRadius
\item OpenRadius
\end{itemize}

Ces trois solutions sont dominantes sur le marché des implémentations Radius open-source.

L'étude destinée à les différencier pour n'en choisir qu'un seul a portée sur plusieurs critères~:

\begin{itemize}
\item Les fonctionnalités supportées.
\item La taille de la communauté et de la documentation.
\item La date de la dernière mise à jour du logiciel.
\end{itemize}

Sur la base de ces trois critères, FreeRadius est ressorti du lot. Outre ses fonctionnalités très complètes, le principal critère différenciateur fut le suivi du projet, avec une dernière mise à jour en 2012, contre 2008 pour GnuRadius et 2007 pour OpenRadius. Le projet FreeRadius bénéficie également d'une communauté exemplaire, et d'une documentation très fournie. Il s'agit d'une référence dans le monde l'entreprise, et assure à B.H. Consulting un suivi du produit sérieux pour l'avenir. C'est également la seule solution qui dispose d'un paquet Debian/Ubuntu, facilitant d'autant plus son installation sur les serveurs.

\subsection{Type de certificats}

%% decrire les differents types en argumentant sur notre choix

\subsection{Interface web}
\subsubsection{Langages et environnements de développement}

Le langage PHP a été imposé pour la conception de l'interface web du projet. Le langage a d'être suffisamement populaire chez les développeurs, pour assurer sa maintenabilité dans le futur. Il permet une approche de la conception propre et cohérente avec les techniques de développement actuelles, en offrant des performances acceptables y compris pour des serveurs de faible capacité.

Une contrainte sur l'environnement de développement est rapidement venue s'ajouter au langage imposé. Les développeurs de B.H. Consulting ont demandé à ce que l'interface exploite les possibilités du \textit{framework} Boostrap Twitter. Ce dernier est particuilèrement populaire depuis sa récente création par l'entreprise du même nom. Il s'agit d'un environnement de développement pour la partie client de l'interface, qui propose une collection très complète de scripts et de feuilles de style pour créer les vues, et offrir une expérience utilisateur confortable. Une formation rapide de cet outil en interne a dû être effectuée par les membres de l'équipe qui ne l'avaient jamais utilisé auparavant.

Du côté du serveur, l'utilisation d'un \textit{framework} est également conseillé, pour forcer une programmation cohérente et organisée de façon stricte. Il existe plusieurs solutions sur le marché, qui ont conduit à une étude au sein de l'équipe, de la meilleure solution à proposer pour la réalisation du projet.

Les plus populaires et puissants qui ont été retenus sont~:

\begin{itemize}
\item CakePHP
\item CodeIgniter
\item Symfony
\item Yii
\item Zend Framework
\end{itemize}

La conclusion de l'étude a poussé l'équipe à se tourner vers la solution CakePHP. Il s'agit en effet d'un des premiers outils de ce type a être développé pour PHP, basé sur la réalisation de Ruby en Rails, le premier du genre. Il bénéficie donc d'une maturité exemplaire, et dispose d'une communauté rassurante, avec une documentation très complète et très détaillée. La prise en main du logiciel est particulièrement rapide, tout en laissant une grande liberté d'adpatation à tous les usages.

CakePHP n'offre pas toutes les possibilités avancées disponibles avec des solutions plus importantes comme Zend Framework ou Symphony, comme les bibliothèques d'interaction avec d'autres services web tiers (comme Google), mais les fonctionnalités fournises suffisent largement pour le projet et les évolutions qu'il pourrait raisonnablement connaitre. Les performances sont particulièrement bonnes en comparaisont à ses concurrents plus complets. La solution Yii aurait permis de construire une appliquation qui tient un peu mieux la charge avec des milliers d'utilisateurs, mais l'usage auquel il est destiné pour le projet n'est pas destiné à prendre en compte de telles problématiques.

CakePHP est disponible sous licence libre (MIT). Les patrons de conception populaires sont disponibles, comme le MVC (modèle-vue-contrôleur) qui permet de segmenter l'application de façon propre et efficace. Un fonctionnement en \textit{ActiveRecord} est également disponible, permettant à des classes PHP d'abstraire totalement la liaison entre l'application et la base de données. Les méthodes du modèle intégrent les opérations CRUD (\textit{create}, \textit{read}, \textit{update}, \textit{delete}) permettant de développer rapidement. De nombreuses autres fonctionnalités sont disponibles, comme un dispacheur d'URL, la validation des données, le moteur de patrons facilitant le formatage, la gestion d'un cache, des composants de sécurité et des scripts de génération de code.

La proposition à l'intervenant industriel de cette solution a particulièrement bien été accueillie puisqu'un développeur de l'entreprise utilise cette solution régulièrement, et sera donc à même de la maintenir dans le futur.

\subsubsection{Gestionnaire de bases de données}

Entre les solutions basées sur des implémentations SQL et les nouvelles possibilité de no-SQL, le choix est vaste pour stocker les informations liées à l'interface. Une solution libre étant demandée par l'entreprise, quelques solutions ne sont dors et déjà pas envisageables. 

La meilleure solution pour faire interagir l'interface web avec le serveur Radius étant de réussir à faire fonctionner ce dernier avec une base de données plutôt qu'une collection de fichiers, le choix de la solution s'est fait en ne gardant que les solutions supportées par le fonctionnement de FreeRadius, l'implémentation de Radius choisie auparavant. Cette contrainte élimine toutes les solutions de no-SQL, et réduit le choix principalement aux solutions populaires MySQL et PostgreSQL, éliminant des solutions plus légères qui auraient pu convenir comme SQLite.

PostgreSQL est beaucoup plus puissant que MySQL, mais il est aussi beaucoup plus lourd et plus contraignant à mettre en place. Les possibilités offertes par MySQL et sa simplicité conviennent parfaitement aux exigences du projet.

Le choix de cette solution a été accueilli favorablement par l'intervenant industriel.

\section{Implémentation}

\end{document}
